% Chapter 6 -- Discussion
% Status: basic draft; revise after the pending results and the advisor's decisions.
\chapter{Discussion}
\label{ch:discussion}

\section{Assessment of the Hypotheses}
\label{sec:disc-hypotheses}

\begin{description}[leftmargin=3em, style=nextline]
  \item[H1 (original): refuted.] The adaptive method does not beat the other placements; it is
        below the private user vector overall and below every method on sparse users
        (Section~\ref{sec:res-main}). \pending{paired tests to state the margins with
        confidence intervals}.
  \item[H2: supported on the mean.] Keeping the user vector private improves \ndcg{} by
        \GainPersonalizedVsBaselinePct\% over averaging it, with everything else equal.
        \pending{paired test}.
  \item[H3: supported.] On the PFedRec model the population-mean head costs
        \DeltaHeadGlobalHead, no $\alpha$ on the grid beats $\alpha=1$, and the data-driven
        rule loses \AbsDeltaHeadAlphaPolicyHc{} (Section~\ref{sec:res-placement}).
  \item[H4: supported; prediction A5 refuted.] Clean post-tuning lowers \ndcg{}
        (\PtDeltaEOneClean, CI \PtCiEOneClean); official-code negatives raise it
        (\PtDeltaEOneLeaky); the contrast is \PtContrastEOne{} and is carried by
        \PtContrastEOneNonZero{} users. The $E=3$ official-code cell exceeded the registered
        ceiling of 0.394 and the published number itself (Section~\ref{sec:res-pfedrec}).
\end{description}

\section{Why Blending with Population Knowledge Did Not Help}
\label{sec:disc-adaptive}

Three mechanisms are consistent with the measurements. First, the population prototype is
nearly uninformative for ranking: the uniform mean of \DsUsers{} heads or vectors cancels
exactly the directions that distinguish one user from another, so any weight $(1-\alpha_u)>0$
on it moves the user toward a point that ranks items like a popularity list. The counterfactual
rows of Table~\ref{tab:head-counterfactuals} show this directly: the population head is worth
less than the own head in every group, and most in the sparse group. Second, the
hierarchical-conditional rule assigns the smallest $\alpha_u$ to exactly those users, so the
rule applies the strongest pull where the prototype is least adequate. Third, the data are not a
cold-start setting: a ``sparse'' MovieLens user still has at least 19 training interactions, enough
to fit a 64-dimensional vector or a 32-dimensional head, so there is little to gain from
borrowing. These observations are consistent with APFL's analysis, in which the optimal mixing
weight approaches the local model as the local and global optima diverge
\cite{deng2020apfl}, and with FedTHE's finding that test-time ensembling of a global and a local
head helps under distribution shift and little in-distribution \cite{jiang2023fedthe}. The
PersonalMLP adds a second problem: \ParamsPersonalMlp{} private parameters trained on a median
of \DsMedianRatings{} interactions in about \ExpectedParticipations{} short visits is a
capacity-to-data ratio under which a non-linear scorer cannot be expected to beat an inner
product. The adaptive design would need a different setting to be useful: a genuine cold-start
population, or a shift between a user's training and test behaviour.

\section{What Post-tuning Measures and What It Does Not}
\label{sec:disc-posttune}

The post-tuning experiment separates three things that the literature reports as one number.
(i) Post-tuning the item table on the client with honest negatives is a \emph{cost} under this
protocol: the client pushes down whatever it draws, and it cannot know that one of the items it
draws is the one the evaluation will ask for. (ii) Excluding the held-out item from the
negatives converts that cost into a gain, because the only rows that can then move among the
candidates are the 99 negatives, which can only move down. This is a property of the evaluation
protocol, not of the method, and it is invisible when the item table is shared, because the
update of one user is diluted among the \ClientsPerRound{} clients of a round. (iii) What
remains of the gap after (ii) is not separated between the participation fraction and the three
deviations K1--K3 of Table~\ref{tab:posttune-deviations}; the official code run on the same
machine (\OfficialCodeNdcg) bounds the sum of those effects if its configuration is confirmed
to be the official default (\pending{confirmation}). The thesis therefore reports PFedRec at
\NdcgPfedrec{} with the shared table as its number, lists the two post-tuning rows next to it,
and states the limitation with the measured cause. \decision{final wording of the limitation
paragraph: (A) two measured causes, $C$ unmeasured but bounded, and the negative-pool leak with
its size; or (B) keep ``no post-tuning at evaluation'' as the cause with a footnote giving the
measurement that has the opposite sign}.

A point on statistics. In the clean cells the mean and the rank test disagree
(Table~\ref{tab:posttune-wilcoxon}): under Pratt's treatment of ties the test almost reduces to
counting the users who gained a little (\PtNPosEOneClean) against those who lost everything
(\PtNNegEOneClean{} at $E=1$, every one of them to \ndcg{} $=0$), and it ignores the size of
the losses. Both views are legitimate answers to different questions; the thesis reports both
and concludes on the mean, because a recommender is judged on its average ranking quality and
a user whose test item disappears from the list is a worse outcome than many users whose item
moves up by one rank.

\section{Fairness of the Placement Comparison}
\label{sec:disc-fairness}

The four cited runs differ in more than the placement (Table~\ref{tab:asymmetries}). PFedRec
uses BCE with four negatives, SGD and $d=\DimPfedrec$; the three others use BPR, Adam and
$d=\DimMf$. Only PFedRec has a final calibration of its private part. The best rounds differ,
and all are selected on the test item. These asymmetries limit the comparison to the reading
``each method at its own, author-faithful configuration'': the placement study is clean only
within the BPR family (Baseline against Personalized), and the PFedRec margin should be read as
a design-level result rather than as a controlled effect of the head alone. \decision{(A) keep
the author-faithful configurations and present Table~\ref{tab:asymmetries} as the reading
guide, or (B) first calibrate the Personalized and Adaptive runs offline to remove the
calibration asymmetry}.

\begin{table}[htbp]
  \caption{Configuration asymmetries between the four cited runs.}
  \label{tab:asymmetries}
  \centering\small
  \fittable{% AUTO-GENERATED by thesis/scripts/build_assets.py -- do not edit by hand.
% configuration asymmetries between the four cited runs
\begin{tabular}{lllll}
\toprule
Aspect & Baseline & Personalized & PFedRec & Adaptive \\
\midrule
Loss & BPR & BPR & BCE (4 negatives) & BPR \\
Embedding dimension & 64 & 64 & 32 & 64 \\
Optimizer & Adam & Adam & SGD (alternating) & Adam \\
Strategy & fedavg & fedavg & fedavg & fedprox \\
Final calibration of the local part & n/a & no & yes (1 epoch, items frozen) & no \\
Best round (selected on the test item) & 83 & 33 & 83 & 83 \\
Extra upload beyond the global part & none & none & none & $p_u$ (prototype) \\
\bottomrule
\end{tabular}
}
\end{table}

\section{Implications for Practice}
\label{sec:disc-practice}

For a cross-device recommender at low participation the results suggest: keep the user-specific
part on the device (it is both the most private and the most accurate choice of the BPR family);
if the model allows it, place personalisation in a private scoring function on a shared item
table rather than in a user vector, and re-align that function to the current table before
serving; do not blend the private part with a population average unless the population is
cold; never upload the user vector; and treat any client-side fine-tuning of the item table at
inference as a protocol decision that must be reported with its negative-sampling rule.
Operationally, evaluate stateful clients with a cache-miss counter, a validity sidecar and a
replay gate, because the failure mode is silent.

\section{Threats to Validity}
\label{sec:disc-threats}

\paragraph{Internal.} One seed per run; no validation item, so best rounds are selected on the
test item (inflation of \InflationBaseline{} to \InflationPfedrec{} points,
Section~\ref{sec:setup-validity}); configuration asymmetries (Section~\ref{sec:disc-fairness});
the adaptive run kept a contrastive term enabled and was not calibrated; the centralised
comparator is untuned. The paired tests, where available, describe users, not seeds.

\paragraph{Construct.} Sampled metrics with 99 negatives are known to disagree with exhaustive
ones \cite{krichene2020sampled}; the comparison is valid because every method ranks the same
candidates, but the absolute numbers are protocol-specific. The ``sparse'' group of MovieLens 1M
is not a cold-start group.

\paragraph{External.} One dataset, one domain (films), one participation fraction. The leak
mechanism of Section~\ref{sec:disc-posttune} applies to any method that fine-tunes a per-user
item table at inference, but its size depends on the number of interactions per user and on
the number of epochs.

\paragraph{Measurement.} All full-population numbers come from runs classified as valid, with
cache-miss counters where the code recorded them; the offline tool reproduced the PFedRec
evaluation bit-exactly. Numbers in this document are generated from the result files by one
script and are not typed by hand.
